\documentclass{article}
\usepackage{subfiles}

\usepackage{geometry}
\usepackage[parfill]{parskip}
\usepackage{mathtools}
\usepackage{tcolorbox}
\usepackage{amsmath,amssymb,amsthm}              
\usepackage{graphicx}
\usepackage{enumitem}
\usepackage{tikz}
\usepackage{mathrsfs}  
\usetikzlibrary{calc}
\usepackage{tikz-cd}
\usepackage{xcolor}
\usepackage{booktabs}
\usepackage[linesnumbered,ruled,vlined]{algorithm2e} 
\usepackage{lineno}
\usepackage[colorlinks]{hyperref}
\usepackage{mdframed}
\usepackage[symbols,nogroupskip,sort=none]{glossaries-extra}
\renewcommand*{\glspostdescription}{\medskip}

\newtheorem*{theorem*}{Theorem}
\newtheorem*{proposition*}{Proposition}
\newtheorem*{corollary*}{Corollary}

\newcommand{\N}{\mathbb{N}}
\newcommand{\Z}{\mathbb{Z}}
\newcommand{\Q}{\mathbb{Q}}
\newcommand{\R}{\mathbb{R}}
\newcommand{\C}{\mathbb{C}}
\newcommand{\G}{\mathbb{G}}
\newcommand{\delayX}{\mathcal{X}}
\newcommand{\X}{\text{x}}
\newcommand{\W}{\text{w}}
\newcommand{\Rho}{\mathrm{P}}
\newcommand{\Lgr}{\mathcal{L}}
\mathchardef\Re="023C
\newcommand{\GLdelay}{\mathrm{GL}_{\mathrm{delay}}}    
\newcommand{\Hrose}{\mathscr{H}} 
\newcommand{\krose}{k_{\theta,\sigma}}


\theoremstyle{plain} 
\newtheorem{theorem}{Theorem}[section]
\newtheorem{lemma}{Lemma}[section]
\newtheorem{corollary}{Corollary}[section]
\newtheorem{proposition}{Proposition}[section]


\theoremstyle{definition} 
\newtheorem{definition}{Definition}[section]

\theoremstyle{remark} 
\newtheorem{remark}{Remark}[section]
\newtheorem{example}{Example}[section]
\newtheorem{assumption}{Assumption}[section]

\DeclareMathOperator{\diag}{diag}
\newcommand{\Var}{\mathrm{Var}}
\newcommand{\Cov}{\mathrm{Cov}}
\DeclareMathOperator{\E}{\mathbb{E}}

% a grey framed box with small padding
\newmdenv[
  linecolor=gray!60,
  linewidth=.8pt,
  topline=true,
  bottomline=true,
  skipabove=\baselineskip,
  skipbelow=\baselineskip,
  innerleftmargin=6pt,
  innerrightmargin=6pt,
  innertopmargin=4pt,
  innerbottommargin=4pt,
  backgroundcolor=gray!5,
  frametitlefont=\normalfont\bfseries,
]{assumptionbox}


% ------------------------------------------------------------------
% A numbered list tailored for assumptions:
%   (W1) ..., (W2) ...   or   (K1) ..., (K2) ...
% ------------------------------------------------------------------
\newlist{assumptionlist}{enumerate}{1}          % base = enumerate, depth = 1
\setlist[assumptionlist]{
    leftmargin=1.4cm,                          % indent
    labelsep=0.2cm,                            % space after label
    itemsep=0.4\baselineskip,                  % vertical gap between items
    label=\textbf{(\Alph*)}                    % default label (over-ridden per list)
}

\pagecolor{white}

\begin{document}

\title{Geometry, Entropy, and Operator Learning with Tunable Kernels}
\author{Patrick S. Mahon}
\maketitle

%%%%%%%%%%%%%%%%%%%%%%%%%%%%%%%%%%%%%%%%%%%%%%%%%%%%%%%%%%%%
% Main paper skeleton
%%%%%%%%%%%%%%%%%%%%%%%%%%%%%%%%%%%%%%%%%%%%%%%%%%%%%%%%%%%%


\tableofcontents


\section{Introduction}
  \subsection{Why tune local kernels?}
    % Drift–diffusion ambiguity in data-driven operators
  \subsection{Key contributions}
    \begin{itemize}[leftmargin=*,nosep]
    \item Closed-form bridge: CGF curvature $\to$ Fisher load $\to$ curvature \& entropy.
    \item Single KL optimisation (little-rose) jointly tunes \textbf{two} dials $(\theta,\sigma)$.
    \item Finite-sample Goldenshluger–Lepski selector with risk guarantees.
    \end{itemize}
  \subsection{Organisation of the paper}

\subfile{body.tex}

\section{Discussion and Outlook}
	\subsection{Key implications}   % ← distilled pay-off list P1–P4
		\begin{enumerate}[label=\textbf{(I\arabic*)},leftmargin=*]
 			\item \textbf{Unified diagnostic of fluctuation strength}\\ The Fisher load $\mathcal I_{\theta,\sigma}$ is \emph{simultaneously} the variance coefficient, the synthetic–curvature prefactor, and the entropy–production rate.  A single dial pair $(\theta, \sigma)$ governs them all.
			
			\item \textbf{Automatic bias–variance balance}\\ 
				Steepest descent of $\mathrm{KL}(\mu_1\parallel k_{\theta,\sigma}\!\#\mu_0)$ is steepest \emph{ascent} of $\mathcal I_{\theta,\sigma}$, so the statistically most informative kernel is found with no cross-validation.
			
			\item \textbf{One-step Schrödinger bridge}\\ 
				The little-rose kernel realises the discrete entropic transport $\mu_0\!\to\!\mu_1$, linking operator learning to optimal-transport theory.
			
			\item \textbf{Landau lens on data}\\
				The Legendre transform of the CGF equips any observable with a Landau-type free-energy landscape whose stiffness is data-computable.
				
			\item \textbf{Lag–dimension selection.}
      				Remark~\ref{rem:delay} suggests using the little‑rose Fisher load as a principled criterion for choosing delay‑embedding parameters.

		\end{enumerate}

	\subsection{Future directions}
  		% brief notes on directed curvature, TDA, entropy cycles, etc.
		
		\begin{enumerate}[label=\textbf{(F\arabic*)},leftmargin=*]
 			
			 \item \textbf{Directed discrete curvature}\\
        				Dial matrices $(\Theta_{ij},\Sigma_{ij})$ define a doubly-weighted digraph, opening analytic work on Ollivier/Bakry–Émery curvature for non-reversible kernels.
				
			\item \textbf{Multiparameter persistent homology}\\
				The two weights furnish a natural bifiltration, inviting multi-parameter TDA of dynamical data sets.
		\end{enumerate}

	\subsection{Limitations}
  		% e.g. local bandwidth assumption, high-dim sparsity

	\subsection{Conclusion}

%%%%%%%%%%%%%%%%%%%%%%%%%%%%%%%%%%%%%%%%%%%%%%%%%%%%%%%%%%%%
% Appendices
%%%%%%%%%%%%%%%%%%%%%%%%%%%%%%%%%%%%%%%%%%%%%%%%%%%%%%%%%%%%
\appendix
\section{Proof details for Sec.~3}
\section{GL selector derivations}
\section{Additional experiments}
\section{Dial-graph topology (exploratory)}

\end{document}
